\section{Related Work}

A compact checklist of rollout-system design choices is provided in \cref{tab:framework_comparison} in the appendix. This section focuses on the qualitative differences behind that comparison.

\subsection{Agent RL Systems}

The first wave of LLM RL infrastructure largely assumed that rollout generation was a Python function owned by the trainer. This assumption is increasingly strained by multi-turn agents, where interaction spans many model calls and environment actions. \prorlagent \citep{prorlagent2026} introduced a service boundary for multi-turn agent rollouts, separating sandbox setup, agent execution, and reward computation from the training process. \polar inherits the same high-level idea that rollout should be a service, but changes the integration contract. Instead of implementing an agent handler inside the rollout service, the user supplies a harness adapter that prepares configuration and launches the native executable. The model proxy then observes the harness from outside.

SkyRL-Agent \citep{skyrlagent2025} is a full-stack system for efficient RL training and evaluation of multi-turn, long-horizon agents, with SkyRL-Gym providing tool-use environments through a Gymnasium-style interface. SkyRL's strength is efficient training once tasks are represented in its environment and agent abstractions. \polar is complementary: it targets the earlier systems problem of running a pre-existing harness whose internal event loop, tool formatting, and context policy should remain unchanged.

PRIME-RL \citep{primerl2026} focuses on large-scale asynchronous RL with trainer-inference separation, stale-policy step semantics, and support for verifiers environments. Slime \citep{slime2025,sglang2024} similarly connects Megatron training with SGLang rollout engines and exposes customizable data-generation interfaces. These systems address the policy-optimization and inference-scaling side of the pipeline. \polar is not a replacement trainer. It is a rollout substrate that can feed asynchronous trainers with trajectories from heavier harnesses than typical verifiable-reward functions.

\subsection{Low-Intrusion Agent Instrumentation}

Agent Lightning \citep{agentlightning2025} proposes a training-agent disaggregation architecture and a unified data interface for converting agent execution into trainable transitions. rLLM \citep{rllm2025} similarly aims to train agents across frameworks with minimal code changes, using tracked clients, decorators, workflow abstractions, and proxy support to collect token IDs and log probabilities. Both systems recognize that researchers should not have to rewrite complete applications to train them.

\polar differs in the chosen minimum integration point. For many coding and terminal agents, the most reliable interface is not an SDK callback graph but the provider API endpoint already used by the harness. The gateway proxy therefore becomes the observation device: it accepts Anthropic, OpenAI Chat, OpenAI Responses, and Google-style requests; translates them to the local inference backend; and records the token-level fields needed by the trainer. This choice is narrower than general observability instrumentation, but it is robust to harnesses implemented as command-line programs, package-managed tools, or binaries.

% OpenClaw-RL is closest in spirit to this proxy-based view. It wraps a self-hosted model behind an OpenAI-compatible API, learns asynchronously from live multi-turn interactions, and targets real-world terminal, GUI, SWE, and tool-call settings \citep{openclawrl2026}. \polar shares the principle that real harness traffic should be a learning source. Its task-submitted rollout service adds a different set of systems choices: explicit runtime isolation, gateway-node scheduling, post-run evaluation, task callbacks, and pluggable trajectory builders/evaluators for benchmark-style RL runs.

\subsection{SWE Task Evaluation and Benchmark}

Harbor \citep{harbor2026} evaluates agents such as Claude Code, OpenHands, Codex CLI, and related systems in containerized environments, supports parallel execution through local and cloud providers, and converts native agent logs into evaluation trajectories. This evaluation-first design is highly aligned with \polar's harness-native motivation. The difference is the model boundary and training data contract. Harbor launches each harness with provider-specific configuration and does not provide a gateway that translates model-provider protocols or mediates the harness's model traffic. As a result, model substitution is limited by what the native harness and external endpoint already support: for example, evaluating a Qwen checkpoint through Claude Code requires an Anthropic-compatible endpoint outside Harbor. \polar instead places a proxy at this boundary, so the same style of harness execution can yield token IDs, log probabilities, loss masks, and rewards that are directly consumable by an RL trainer.

SWE-bench \citep{swebench2024} established real GitHub issue resolution as a benchmark requiring repository understanding, editing, and executable validation. SWE-Gym \citep{swegym2024} extends this direction with training environments, verifiers, and trajectories for software-engineering agents. These workloads are a natural stress test for rollout infrastructure because they combine expensive runtime setup, sparse patch-level rewards, long-tail execution time, and many opportunities for harness-side state to diverge from a clean evaluator state.

\subsection{Token Fidelity and Retokenization Drift}

The training signal in agent RL is only correct if it is attached to the tokens sampled by the behavior policy. This is difficult in agent harnesses because provider APIs may return text, tool-call JSON, reasoning fields, or streamed events rather than the exact token IDs and log probabilities used by the inference backend. The vLLM and Agent Lightning discussion of retokenization drift emphasizes that decoding and re-encoding a transcript can produce different token IDs from the original generation \citep{vllmretokenization2025}. \polar follows the same token-fidelity principle but applies it to arbitrary harness rollouts: generated assistant tokens are copied from inference responses, non-generated interstitial tokens are taken from canonical prompt tokenization, and the loss mask marks only behavior-policy tokens as trainable.
